\documentclass[../../main.tex]{subfiles}
\begin{document}
\begin{oframed}
  \subfile{problem}
\end{oframed}

{\bf [解答]}

$e(\theta)=\cos\theta+i\sin\theta$とし，複素数平面で置き換えて考える．
$\overrightarrow{P_{n-1}P_n}$を表す複素数を$a_n$とすれば，題意より
\begin{align}
  \left\{
    \begin{array}{l}
      a_{n+1}=2e\left(\dfrac{2\pi}{3}\right)a_n  \\
      a_1=1
    \end{array}
    \right.
  \end{align}
  だから，繰り返し用いて$a_n=\left\{2e\left(\dfrac{2\pi}{3}\right)\right\}^{n-1}$である．
  従って$P_{3n}$を表す複素数$b_{n}$として
  $a=2e\left(\dfrac{2\pi}{3}\right)$とおけば，
  \begin{align}
    b_n&=\sum_{k=1}^{3n}(a_k) \\
    &=\sum_{k=1}^{3n}(a^{k-1}) \\
    &=\dfrac{1-a^{3n}}{1-a} \\
    &=\dfrac{2+\sqrt{3}i}{7}(1-2^{3n})
  \end{align}
  であるから，求める座標は
  \begin{align}
    P_{3n}\left(\dfrac{2}{7}(1-2^{3n}),\dfrac{\sqrt{3}}{7}(1-2^{3n})\right)
  \end{align}
  である．$\cdots$(答)

  \newpage
  \end{document}
